\documentclass[aspectratio=169,11pt]{beamer}

\usetheme{Madrid}
\usecolortheme{default}

\usepackage[T1]{fontenc}
\usepackage[utf8]{inputenc}
\usepackage{amsmath,amssymb}
\usepackage{booktabs}
\usepackage{listings}
\usepackage{xcolor}

\definecolor{LeanKeyword}{RGB}{0,76,153}
\definecolor{LeanComment}{RGB}{90,110,90}
\definecolor{LeanString}{RGB}{150,70,40}
\definecolor{CodeBg}{RGB}{248,248,248}

\lstdefinelanguage{Lean}{
  morekeywords={
    import,namespace,end,section,variable,universe,inductive,structure,class,
    instance,where,def,theorem,example,by,match,with,fun,let,if,then,else,
    deriving,Type,Prop
  },
  sensitive=true,
  morecomment=[l]--,
  morestring=[b]"
}

\lstdefinestyle{leanstyle}{
  language=Lean,
  basicstyle=\ttfamily\small,
  keywordstyle=\color{LeanKeyword}\bfseries,
  commentstyle=\color{LeanComment},
  stringstyle=\color{LeanString},
  backgroundcolor=\color{CodeBg},
  showstringspaces=false,
  columns=fullflexible,
  keepspaces=true,
  frame=single,
  framerule=0.2pt,
  xleftmargin=0.5em,
  xrightmargin=0.5em
}

\lstset{style=leanstyle}

\newcommand{\DTwoPartName}{Overview}
\newcommand{\setDTwoPart}[1]{\gdef\DTwoPartName{#1}}
\newcommand{\leanfile}[1]{\vfill{\small\color{gray}Live Lean: \texttt{#1}}}
\newcommand{\partlink}[2]{\hyperlink{#1}{\textcolor{LeanKeyword}{\textbf{#2}}}}

\setbeamertemplate{footline}{
  \leavevmode%
  \hbox{%
    \begin{beamercolorbox}[wd=.30\paperwidth,ht=2.6ex,dp=1ex,leftskip=1em]{author in head/foot}%
      \usebeamerfont{author in head/foot}\DTwoPartName
    \end{beamercolorbox}%
    \begin{beamercolorbox}[wd=.50\paperwidth,ht=2.6ex,dp=1ex,center]{title in head/foot}%
      \usebeamerfont{title in head/foot}\insertshorttitle
    \end{beamercolorbox}%
    \begin{beamercolorbox}[wd=.20\paperwidth,ht=2.6ex,dp=1ex,rightskip=1em plus1fil]{date in head/foot}%
      \usebeamerfont{date in head/foot}\insertframenumber{} / \inserttotalframenumber
    \end{beamercolorbox}%
  }%
  \vskip0pt%
}

\title[D2: Inductive Types and Classes]{D2: Inductive Types and Classes}
\subtitle{Constructors, recursion, induction, structures, and type classes}
\author{Summer School on Lean}
\date{July 14, Tuesday}

\begin{document}

\begin{frame}
  \titlepage
\end{frame}

\begin{frame}{Learning Goals}
  

  \begin{itemize}
    \item read simple \texttt{inductive}, \texttt{structure}, and \texttt{class} declarations;
    \item predict the cases of a pattern match from the constructors of a type;
    \item recognize the shape of recursive definitions and induction proofs;
    \item use basic examples of \texttt{Option}, \texttt{List}, \texttt{Fin}, and \texttt{Multiset};
    \item understand what \texttt{instance} search and \texttt{deriving} are for.
  \end{itemize}
\end{frame}

\begin{frame}{Four-Part Outline}
  \begin{tabular}{@{}p{0.18\textwidth}p{0.72\textwidth}@{}}
    \toprule
    Part I & \partlink{part-inductive}{Inductive} \\
           & Motivation and basic \texttt{inductive} declarations \\
    \midrule
    Part II & \partlink{part-recursion-induction}{Recursion and Induction} \\
            & Pattern matching, recursion, \texttt{cases}, and \texttt{induction} \\
    \midrule
    Part III & \partlink{part-structures-apis}{Structures and APIs} \\
             & \texttt{structure}, \texttt{Option}, \texttt{List}, \texttt{Fin}, \texttt{Multiset} \\
    \midrule
    Part IV & \partlink{part-type-classes}{Type Classes} \\
            & \texttt{class}, \texttt{instance}, \texttt{deriving}, and an integrated expression example \\
    \bottomrule
  \end{tabular}
\end{frame}

\begin{frame}{Timing}
  \begin{tabular}{@{}ll@{}}
    \toprule
    Time & Topic \\
    \midrule
    0--35 min & Part I: Inductive \\
    35--80 min & Part II: Recursion and Induction \\
    80--100 min & Break \\
    100--145 min & Part III: Structures and APIs \\
    145--180 min & Part IV: Type Classes \\
    \bottomrule
  \end{tabular}
\end{frame}

\setDTwoPart{Part I: Inductive}
\section{Motivation}

\begin{frame}{Data Built from Constructors}
  \hypertarget{part-inductive}{}
  Many Lean types are described by listing their constructors:

  \begin{itemize}
    \item \texttt{Bool}: \texttt{true}, \texttt{false}
    \item \texttt{Nat}: \texttt{zero}, \texttt{succ}
    \item \texttt{Option alpha}: \texttt{none}, \texttt{some a}
    \item \texttt{List alpha}: \texttt{[]}, \texttt{x :: xs}
  \end{itemize}

  Knowing how values are built tells us how to compute with them and how to prove statements about them.
\end{frame}

\begin{frame}[fragile]{Weekday}
  \begin{lstlisting}
inductive Weekday where
  | mon | tue | wed | thu | fri | sat | sun
  deriving Repr, BEq, DecidableEq
  \end{lstlisting}

  \begin{itemize}
    \item \texttt{Weekday} is a new type with seven constructors.
    \item The constructors are the only ways to build a \texttt{Weekday}.
    \item \texttt{deriving} generates printing and equality automatically.
  \end{itemize}
\end{frame}

\section{Inductive Types}

\begin{frame}{Reading an Inductive Declaration}
  When you encounter an \texttt{inductive} declaration, ask:

  \begin{enumerate}
    \item What is the new type?
    \item What are the constructors?
    \item Which constructors carry data?
    \item Which constructors are recursive?
  \end{enumerate}

  These four answers predict the shape of pattern matches, recursive definitions, and induction proofs over the type.

  \leanfile{code/D2\_Inductive.lean}
\end{frame}

\begin{frame}[fragile]{Constructors Are Functions}
  \begin{lstlisting}
inductive MyNat where
  | zero : MyNat
  | succ : MyNat -> MyNat
  \end{lstlisting}

  \begin{itemize}
    \item \texttt{zero} is already a natural number.
    \item \texttt{succ} needs a smaller natural number before it produces a new one.
    \item Use \texttt{\#check} to inspect constructor types.
  \end{itemize}
\end{frame}

\begin{frame}[fragile]{Parameters vs. Constructor Arguments}
  \begin{lstlisting}
inductive MyOption (alpha : Type) where
  | none : MyOption alpha
  | some : alpha -> MyOption alpha
  \end{lstlisting}

  \begin{itemize}
    \item \texttt{alpha} is a parameter of the whole family of types.
    \item \texttt{some} stores one particular value of type \texttt{alpha}.
    \item Examples: \texttt{MyOption Nat}, \texttt{MyOption String}.
  \end{itemize}
\end{frame}

\begin{frame}[fragile]{Recursive Constructors}
  \begin{lstlisting}
inductive MyNat where
  | zero : MyNat
  | succ : MyNat -> MyNat

inductive MyList (alpha : Type) where
  | nil : MyList alpha
  | cons : alpha -> MyList alpha -> MyList alpha
  \end{lstlisting}

  \begin{itemize}
    \item Recursive constructors mention the type being defined.
    \item The recursive argument is the smaller object.
    \item This is why structural recursion and induction follow the same shape.
  \end{itemize}
\end{frame}

\begin{frame}[fragile]{Addition on MyNat}
  \begin{lstlisting}
def MyNat.add : MyNat -> MyNat -> MyNat
  | m, zero => m
  | m, succ n => succ (MyNat.add m n)
  \end{lstlisting}

  \begin{itemize}
    \item The two equations follow the two constructors of \texttt{MyNat}.
    \item Recursion is on the second argument: \texttt{succ n} reduces to \texttt{n}.
    \item These are the familiar recursive equations \(m + 0 = m\) and \(m + (n+1) = (m+n)+1\).
  \end{itemize}
\end{frame}

\begin{frame}[fragile]{Tree}
  \begin{lstlisting}
inductive Tree (alpha : Type) where
  | leaf : alpha -> Tree alpha
  | node : Tree alpha -> Tree alpha -> Tree alpha
  \end{lstlisting}

  \begin{itemize}
    \item A leaf stores one value.
    \item A node stores two subtrees of the same type.
    \item In the node case of an induction proof there are two induction hypotheses, one per subtree.
  \end{itemize}
\end{frame}

\setDTwoPart{Part II: Recursion and Induction}
\section{Pattern Matching and Recursion}

\begin{frame}{Part II: Recursion and Induction}
  \hypertarget{part-recursion-induction}{}
  Constructors are used in two more ways:

  \begin{itemize}
    \item pattern matching to inspect a value of an inductive type;
    \item structural recursion to define functions over recursive data.
  \end{itemize}

  Induction proofs then follow the same case structure as these definitions.
\end{frame}

\begin{frame}[fragile]{Pattern Matching}
  \begin{lstlisting}
def isWeekend : Weekday -> Bool
  | Weekday.sat => true
  | Weekday.sun => true
  | _ => false
  \end{lstlisting}

  \begin{itemize}
    \item Pattern matching eliminates an inductive value.
    \item Each branch handles one possible constructor shape.
    \item The wildcard pattern \texttt{\_} handles remaining cases.
  \end{itemize}

  \leanfile{code/D2\_PatternMatching.lean}
\end{frame}

\begin{frame}[fragile]{Matching on Option}
  \begin{lstlisting}
def getWithDefault {alpha : Type}
    (default : alpha) : Option alpha -> alpha
  | none => default
  | some a => a
  \end{lstlisting}

  \begin{itemize}
    \item The \texttt{none} branch has no stored value.
    \item The \texttt{some} branch exposes the stored value.
    \item The return type is the same in both branches.
  \end{itemize}
\end{frame}

\begin{frame}[fragile]{Recursive Definitions on Natural Numbers}
  \begin{lstlisting}
def add (m : Nat) : Nat -> Nat
  | 0 => m
  | n + 1 => add m n + 1

def mul (m : Nat) : Nat -> Nat
  | 0 => 0
  | n + 1 => mul m n + m
  \end{lstlisting}

  \begin{itemize}
    \item The recursion is on the second argument.
    \item The equations are ordinary recursive mathematical definitions.
    \item Powers can be defined in the same style.
  \end{itemize}
\end{frame}

\begin{frame}[fragile]{Recursive Definitions on Lists}
  \begin{lstlisting}
def myLength {alpha : Type} : List alpha -> Nat
  | [] => 0
  | _ :: xs => myLength xs + 1
  \end{lstlisting}

  \begin{itemize}
    \item The recursive call is made on \texttt{xs}, a smaller list.
    \item Lean accepts this as structural recursion.
    \item The two branches match the two list constructors.
  \end{itemize}
\end{frame}

\begin{frame}[fragile]{Map on Lists}
  \begin{lstlisting}
def myMap {alpha beta : Type} (f : alpha -> beta) :
    List alpha -> List beta
  | [] => []
  | x :: xs => f x :: myMap f xs
  \end{lstlisting}

  \begin{itemize}
    \item The empty list maps to the empty list.
    \item The \texttt{cons} branch applies \texttt{f} to the head and recurses on the tail.
    \item The output list has the same length as the input list.
  \end{itemize}
\end{frame}

\section{Induction Proofs}

\begin{frame}{From Recursion to Induction}
  \begin{tabular}{@{}lll@{}}
    \toprule
    Data & Program & Proof \\
    \midrule
    base constructor & base branch & base case \\
    recursive constructor & recursive call & induction hypothesis \\
    \bottomrule
  \end{tabular}

  \vspace{0.8em}

  A function definition on an inductive type and an induction proof over that type share the same case structure.
\end{frame}

\begin{frame}[fragile]{Proof by Cases}
  \begin{lstlisting}
example (b : Bool) : b = true \/ b = false := by
  cases b <;> simp
  \end{lstlisting}

  \begin{itemize}
    \item \texttt{cases} splits a value according to its constructors.
    \item It is best for small or non-recursive inductive types.
  \end{itemize}

  \leanfile{code/D2\_InductionProofs.lean}
\end{frame}

\begin{frame}[fragile]{Induction on Lists}
  \begin{lstlisting}
theorem append_nil_right {alpha : Type} (xs : List alpha) :
    xs ++ [] = xs := by
  induction xs with
  | nil =>
      rfl
  | cons x xs ih =>
      rw [ih]
  \end{lstlisting}

  \begin{itemize}
    \item The base case is the empty list.
    \item The step case is \texttt{x :: xs}.
    \item The induction hypothesis is the theorem already known for \texttt{xs}.
  \end{itemize}
\end{frame}

\begin{frame}[fragile]{Induction on Addition}
  \begin{lstlisting}
theorem zero_add (n : Nat) : add 0 n = n := by
  induction n with
  | zero =>
      rfl
  | succ n ih =>
      calc
        add 0 (n + 1) = add 0 n + 1 := rfl
        _ = n + 1 := by rw [ih]
  \end{lstlisting}

  \begin{itemize}
    \item The base case proves the statement for \texttt{0}.
    \item The step case uses the induction hypothesis for \texttt{n}.
    \item This is the familiar proof of a recursive equation by induction.
  \end{itemize}
\end{frame}

\setDTwoPart{Break}
\section{Break}

\begin{frame}{Break}
  \begin{center}
    \Huge 20-minute break
  \end{center}

  

  
\end{frame}

\setDTwoPart{Part III: Structures and APIs}
\section{Structures}

\begin{frame}{Part III: Structures and APIs}
  \hypertarget{part-structures-apis}{}
  Three topics:

  \begin{itemize}
    \item \texttt{structure}: an inductive type with one constructor and named fields;
    \item the standard inductive types \texttt{Option}, \texttt{List}, \texttt{Fin}, \texttt{Multiset};
    \item how each one supports a different modeling choice.
  \end{itemize}
\end{frame}

\begin{frame}[fragile]{Structures as Records}
  \begin{lstlisting}
structure Vector2 where
  x : Int
  y : Int
  \end{lstlisting}

  \begin{itemize}
    \item A structure bundles named fields.
    \item Values can be constructed with record syntax.
    \item Fields can be accessed with dot notation.
  \end{itemize}

  \leanfile{code/D2\_Structure.lean}
\end{frame}

\begin{frame}[fragile]{Operations on a Structure}
  \begin{lstlisting}
def Vector2.add (u v : Vector2) : Vector2 :=
  { x := u.x + v.x, y := u.y + v.y }

def Vector2.dot (u v : Vector2) : Int :=
  u.x * v.x + u.y * v.y
  \end{lstlisting}

  \begin{itemize}
    \item \texttt{u.x} and \texttt{u.y} access coordinates.
    \item Record syntax constructs new vectors from field values.
    \item This mirrors ordinary coordinate definitions in linear algebra.
  \end{itemize}
\end{frame}

\begin{frame}[fragile]{Bundling Data with a Property}
  \begin{lstlisting}
structure PointOnDiagonal where
  x : Int
  y : Int
  onDiagonal : x = y
  \end{lstlisting}

  \begin{itemize}
    \item Structures can contain data fields and proof fields.
    \item The field \texttt{onDiagonal} records a mathematical property.
    \item This pattern is common in formalized mathematics.
  \end{itemize}
\end{frame}

\begin{frame}{Inductive vs. Structure}
  \begin{tabular}{@{}lll@{}}
    \toprule
    Tool & Typical use & Shape \\
    \midrule
    \texttt{inductive} & alternatives, recursive data & one or more constructors \\
    \texttt{structure} & bundled fields & one constructor with named fields \\
    \bottomrule
  \end{tabular}

  \vspace{1em}

  Many mathematical objects in Lean are represented as structures.
\end{frame}

\section{Common Inductive Types}

\begin{frame}{Common Inductive Types}
  \begin{tabular}{@{}lll@{}}
    \toprule
    Type & Meaning & Typical use \\
    \midrule
    \texttt{Option alpha} & a value or nothing & partial lookup \\
    \texttt{List alpha} & finite ordered data & recursion and iteration \\
    \texttt{Fin n} & a natural number below \texttt{n} & safe indexing \\
    \texttt{Multiset alpha} & unordered data with multiplicity & finite combinatorics \\
    \bottomrule
  \end{tabular}

  \leanfile{code/D2\_CommonTypes.lean}
\end{frame}

\begin{frame}[fragile]{Option}
  \begin{lstlisting}
def safeHead {alpha : Type} : List alpha -> Option alpha
  | [] => none
  | x :: _ => some x
  \end{lstlisting}

  \begin{itemize}
    \item The return type signals that the lookup may have no result.
    \item Any caller has to handle both \texttt{none} and \texttt{some}.
    \item Partiality is recorded in the type, not in a side channel.
  \end{itemize}
\end{frame}

\begin{frame}{List}
  Common library operations on \texttt{List alpha}:

  \begin{itemize}
    \item \texttt{map}: transform each element;
    \item \texttt{filter}: keep elements satisfying a predicate;
    \item \texttt{foldl}, \texttt{foldr}: accumulate values from left or right;
    \item \texttt{length}, \texttt{++} (append), \texttt{mem}.
  \end{itemize}

  Each operation packages a standard recursive pattern over \texttt{nil} and \texttt{cons}.
\end{frame}

\begin{frame}[fragile]{Fin n}
  \begin{lstlisting}
#check Fin 3
#check (0 : Fin 3)
#check (2 : Fin 3)
  \end{lstlisting}

  \begin{itemize}
    \item An element of \texttt{Fin n} is a natural number together with a proof that it is less than \texttt{n}.
    \item Writing \texttt{(3 : Fin 3)} is a type error.
    \item \texttt{Fin n} is the standard index type for vectors, matrices, and finite choices.
  \end{itemize}
\end{frame}

\begin{frame}{Multiset}
  \begin{itemize}
    \item A \texttt{List alpha} remembers order; \texttt{[1, 2, 1]} and \texttt{[1, 1, 2]} are different lists.
    \item A \texttt{Multiset alpha} forgets order but remembers multiplicity; the same two examples are equal as multisets.
    \item \texttt{Multiset} is defined in mathlib using \texttt{List} quotiented by permutation.
  \end{itemize}
\end{frame}

\setDTwoPart{Part IV: Type Classes}
\section{Type Classes}

\begin{frame}{From Structures to Classes}
  \hypertarget{part-type-classes}{}
  \begin{itemize}
    \item A \texttt{structure} bundles data and operations into one type.
    \item A \texttt{class} is a structure whose instances Lean searches for automatically.
    \item A parameter written as \texttt{[C alpha]} asks Lean to supply a \texttt{C alpha} instance at every call site.
    \item This is how \texttt{+}, \texttt{*}, and algebraic structure work uniformly across many types.
  \end{itemize}
\end{frame}

\begin{frame}[fragile]{A Small Type Class}
  \begin{lstlisting}
class HasNorm (alpha : Type) where
  norm : alpha -> Nat

def doubleNorm [HasNorm alpha] (x : alpha) : Nat :=
  2 * HasNorm.norm x
  \end{lstlisting}

  \begin{itemize}
    \item \texttt{[HasNorm alpha]} is a type class parameter.
    \item Lean tries to synthesize the instance automatically.
  \end{itemize}

  \leanfile{code/D2\_Typeclass.lean}
\end{frame}

\begin{frame}[fragile]{Instances}
  \begin{lstlisting}
structure Vector2 where
  x : Int
  y : Int

instance : HasNorm Vector2 where
  norm v := v.x.natAbs + v.y.natAbs
  \end{lstlisting}

  \begin{itemize}
    \item An \texttt{instance} provides the data a class requires for a particular type.
    \item Once registered, the instance is found automatically wherever \texttt{[HasNorm Vector2]} is required.
    \item \texttt{Int.natAbs} sends an integer to its absolute value as a \texttt{Nat}.
    \item Use \texttt{\#synth HasNorm Vector2} to ask Lean what it found.
  \end{itemize}
\end{frame}

\begin{frame}[fragile]{Add and Mul as Classes}
  \begin{lstlisting}
#synth Add Nat
#synth Add Int
#synth Mul Nat
#synth Mul Int

def addThree [Add alpha] (a b c : alpha) : alpha :=
  a + b + c
  \end{lstlisting}

  \begin{itemize}
    \item \texttt{+} is the method of the class \texttt{Add}; \texttt{*} is the method of \texttt{Mul}.
    \item \texttt{addThree} works for any \texttt{alpha} with an \texttt{Add} instance.
    \item Built-in instances exist for \texttt{Nat} and \texttt{Int}; mathlib adds many more.
  \end{itemize}
\end{frame}

\begin{frame}[fragile]{Mathlib: Algebraic Instances}
  \begin{lstlisting}
import Mathlib

#synth Add Nat
#synth Add Int
#synth Add Rat
#synth Add Complex

#synth Ring Int
#synth Field Rat
#synth Field Complex
  \end{lstlisting}

  \begin{itemize}
    \item Mathlib provides a rich algebraic hierarchy.
    \item Lean can find instances for familiar number systems.
    \item This is why generic algebraic theorems can apply broadly.
  \end{itemize}

  \leanfile{code/D2\_Typeclass\_Mathlib.lean}
\end{frame}

\begin{frame}{Classes for Mathematical Structure}
  \begin{tabular}{@{}ll@{}}
    \toprule
    Class & Purpose \\
    \midrule
    \texttt{Add}, \texttt{Mul} & notation for operations \\
    \texttt{Zero}, \texttt{One} & distinguished elements \\
    \texttt{Semigroup}, \texttt{Monoid} & associative operations, identity \\
    \texttt{Group} & inverse operation \\
    \texttt{Ring}, \texttt{CommRing} & ring structure \\
    \texttt{Field} & field structure \\
    \bottomrule
  \end{tabular}
\end{frame}

\begin{frame}{Operations vs.\ Laws}
  \begin{itemize}
    \item \texttt{[Add alpha]} supplies only the operation \texttt{+ : alpha -> alpha -> alpha}.
    \item It does not by itself give associativity, commutativity, or an identity.
    \item For those laws, use stronger classes: \texttt{AddCommMonoid alpha}, \texttt{Ring alpha}, \texttt{Field alpha}.
    \item A theorem stated with \texttt{[Ring R]} applies to every type for which Lean can find a \texttt{Ring} instance.
  \end{itemize}
\end{frame}

\begin{frame}[fragile]{Deriving}
  \begin{lstlisting}
inductive Color where
  | red | green | blue
  deriving Repr, BEq, DecidableEq
  \end{lstlisting}

  \begin{itemize}
    \item \texttt{deriving} generates standard instances.
    \item It makes custom types usable in small examples immediately.
  \end{itemize}
\end{frame}

\section{Integrated Example}

\begin{frame}[fragile]{Expression Trees}
  \begin{lstlisting}
inductive Expr where
  | const : Nat -> Expr
  | add : Expr -> Expr -> Expr
  | var : String -> Expr
  deriving Repr, BEq, DecidableEq
  \end{lstlisting}

  Syntax trees are naturally inductive.

  \leanfile{code/D2\_IntegratedExample.lean}
\end{frame}

\begin{frame}{Exercises}
  \begin{itemize}
    \item Complete the guided holes in \texttt{D2\_Exercises.lean}.
    \item Add a \texttt{mul : Expr -> Expr -> Expr} constructor to \texttt{Expr}.
    \item Extend \texttt{eval} and \texttt{vars} to the new constructor.
    \item Prove \texttt{nodeCount (mirror e) = nodeCount e} by induction on \texttt{e}.
  \end{itemize}

  \leanfile{code/D2\_Exercises.lean}
\end{frame}

\end{document}
